\subsection{Factorialidad de los anillos locales regulares}
\label{subsection:IV.21.11}
\label{IV.21.11}

\begin{theorem}[21.11.1]
\label{IV.21.11.1}
\textup{(Auslander-Buchsbaum)} ---
Un anillo local noetheriano regular es factorial.
\end{theorem}

\begin{proof}
La demostración que sigue se debe a I.~Kaplansky.

Sea $A$ un anillo local noetheriano regular de dimensión $n$; razonaremos
por inducción sobre $n$. Para $n=0$, $A$ es un cuerpo y, para $n=1$, un anillo de
valoración discreta, por tanto principal (y \emph{a fortiori} factorial). Supongamos, pues,
$n\geq2$ y demostrado el teorema para los anillos regulares de dimensión $<n$.
Pongamos $X=\operatorname{Spec}(A)$, denotemos por $a$ el punto cerrado de $X$ y pongamos
$U=X-\{a\}$. En todo punto $y\in U$, se tiene
$\dim(\mathcal O_{X,y})\leq n-1$ y, puesto que $A$ es regular, los anillos
$\mathcal O_{X,y}$ también lo son $(\mathbf 0, 17.3.2)$; por tanto, la hipótesis de
inducción implica que son factoriales. Además, se tiene
$\operatorname{prof}(A)=\dim(A)\geq2$, ya que $A$ es regular y, por tanto, de
Cohen--Macaulay $(\mathbf 0, 17.1.3)$.
\oldpage[IV-4]{303}
Utilizando (21.6.14), nos reducimos a demostrar que $\operatorname{Pic}(U)=0$.
Consideremos, pues, un $\mathcal O_U$-módulo invertible $\mathscr L$, y demostremos que es
isomorfo a $\mathcal O_U$. Resulta de $(\mathbf I, 9.4.5)$ que existe un
$\mathcal O_X$-módulo \emph{coherente} $\mathscr F$ tal que
$\mathscr F|U=\mathscr L$. Como $A$ es regular, y por tanto de dimensión cohomológica
finita $(\mathbf 0, 17.3.1)$, existe una resolución libre por la izquierda \emph{finita} de
$\mathscr F$:
\[
  0\leftarrow\mathscr F\leftarrow\mathcal O_X^{n_1}
  \leftarrow\mathcal O_X^{n_2}\leftarrow\cdots
  \leftarrow\mathcal O_X^{n_h}\leftarrow0
\]
$(\mathbf 0, 17.2.8$ y $\mathbf 0, 17.2.2$, (iii)). Por restricción a $U$, se tiene
entonces una resolución finita
\[
\label{IV.21.11.1.1}
  0\leftarrow\mathscr L\leftarrow\mathcal O_U^{n_1}
  \leftarrow\mathcal O_U^{n_2}\leftarrow\cdots
  \leftarrow\mathcal O_U^{n_h}\leftarrow0.
  \tag{21.11.1.1}
\]
\end{proof}

El teorema resultará entonces de las consideraciones generales siguientes. Sea,
sobre un espacio anillado $X$, $\mathscr E$ un $\mathcal O_X$-módulo localmente libre de rango finito;
denotaremos por $\Lambda^{\max}\mathscr E$ el $\mathcal O_X$-módulo invertible que,
en un entorno de cada punto de $X$, es igual a $\Lambda^p\mathscr E$, donde
$p$ denota el rango de $\mathscr E$ en este entorno (que puede variar con la componente
conexa de $X$). Con esta notación, se tiene el

\begin{lemma}[21.11.1.2]
\label{IV.21.11.1.2}
Sean $X$ un espacio anillado en anillos locales, y
\[
  0\leftarrow\mathscr E_0\xleftarrow{u_0}\mathscr E_1
  \leftarrow\cdots\leftarrow\mathscr E_h\leftarrow0
\]
una sucesión exacta de $\mathcal O_X$-módulos localmente libres de rango finito; entonces el
$\mathcal O_X$-módulo invertible
\[
  \bigotimes_{0\leq i\leq h}
  (\Lambda^{\max}\mathscr E_i)^{\otimes(-1)^i}
\]
es isomorfo a $\mathcal O_X$.
\end{lemma}

\begin{proof}
Demostremos cómo este lema permitirá concluir la demostración de (21.11.1). Para ello basta
observar que, para todo entero $n$,
$\Lambda^{\max}(\mathcal O_U^n)=\Lambda^n\mathcal O_U^n$ es isomorfo a
$\mathcal O_U$, al igual que $\mathcal O_U^{\otimes(-1)}$. Como, por otra parte,
$\Lambda^{\max}\mathscr L=\mathscr L$ para todo $\mathcal O_U$-módulo invertible
$\mathscr L$, el lema (21.11.1.2), aplicado a la sucesión exacta (21.11.1.1), muestra
que $\mathscr L$ es isomorfo a $\mathcal O_U$.

Queda por demostrar (21.11.1.2); se procede por inducción sobre $h$, siendo trivial el lema
para $h=1$. Para $h>1$, $\mathscr N=\operatorname{Ker}(u_0)$ es un
$\mathcal O_X$-módulo localmente libre de rango finito $(\mathbf 0_{\mathrm I}, 5.5.5)$,
y se tienen las dos sucesiones exactas
\[
  0\leftarrow\mathscr E_0\leftarrow\mathscr E_1
  \leftarrow\mathscr N\leftarrow0
\]
\[
  0\leftarrow\mathscr N\leftarrow\mathscr E_2
  \leftarrow\cdots\leftarrow\mathscr E_h\leftarrow0.
\]
En virtud de la hipótesis de inducción,
$(\Lambda^{\max}\mathscr N)\otimes
(\bigotimes_{2\leq i\leq h}(\Lambda^{\max}\mathscr E_i)^{\otimes(-1)^{i-1}})$
es isomorfo a $\mathcal O_X$; por tanto, bastará definir un isomorfismo canónico
\[
  (\Lambda^{\max}\mathscr N)\otimes(\Lambda^{\max}\mathscr E_0)
  \xrightarrow{\sim}\Lambda^{\max}\mathscr E_1
\]
para terminar la demostración. Ahora bien, existe un recubrimiento abierto $(U_\alpha)$ de
$X$ tal que, en todo $U_\alpha$, $\mathscr E_1|U_\alpha$ sea suma directa de
$\mathscr N|U_\alpha$ y de un $\mathcal O_{U_\alpha}$-módulo localmente libre
$\mathscr M_\alpha$ $(\mathbf 0_{\mathrm I}, 5.5.5)$, de donde resulta un isomorfismo canónico
$v_\alpha:\mathscr M_\alpha\xrightarrow{\sim}\mathscr E_0|U_\alpha$.
Como se tiene un isomorfismo canónico
\[
  r_\alpha:(\Lambda^{\max}\mathscr N|U_\alpha)
  \otimes(\Lambda^{\max}\mathscr M_\alpha)
  \xrightarrow{\sim}(\Lambda^{\max}\mathscr E_1)|U_\alpha,
\]
se deduce de él mediante $v_\alpha$ un isomorfismo
\[
  u_\alpha:(\Lambda^{\max}\mathscr N|U_\alpha)
  \otimes(\Lambda^{\max}\mathscr E_0|U_\alpha)
  \xrightarrow{\sim}\Lambda^{\max}\mathscr E_1|U_\alpha.
\]
\oldpage[IV-4]{304}
Se trata de demostrar que $u_\alpha$ y $u_\beta$ coinciden en
$U_\alpha\cap U_\beta$ para dos índices cualesquiera $\alpha$, $\beta$. Ahora bien, si
$v'_\alpha$ y $v'_\beta$ son las restricciones a $U_\alpha\cap U_\beta$ de
$v_\alpha$ y $v_\beta$ respectivamente, se tiene
$v'_\alpha=v'_\beta\circ w_{\beta\alpha}$, donde
\[
  w_{\beta\alpha}:\mathscr M_\alpha|(U_\alpha\cap U_\beta)
  \xrightarrow{\sim}\mathscr M_\beta|(U_\alpha\cap U_\beta)
\]
es la «proyección paralela a $\mathscr N$» tal que, para toda sección
$s\in\Gamma(U_\alpha\cap U_\beta,\mathscr M_\alpha)$,
$w_{\beta\alpha}(s)=s+t$, con
$t\in\Gamma(U_\alpha\cap U_\beta,\mathscr N)$; la identidad de $u_\alpha$ y de
$u_\beta$ resulta inmediatamente de este hecho y de la definición del isomorfismo canónico
$r_\alpha$ (Bourbaki, \emph{Alg.}, cap.~III, 3\textsuperscript{a} ed.).
\end{proof}
